% !TEX root = ../main.tex
% capitulos/4_consideracoes_finais.tex

Ao final do texto e antes de apresentar as referências é obrigatório apresentar o quadro com informações adicionais, conforme descrito abaixo.

Apresentar agradecimentos a todas pessoas e entidades que participaram de forma direta ou indireta para a construção da pesquisa e do artigo, agências de fomento, autorizações de pesquisa, além informar as possíveis agências de fomento e licenças concedidas para a realização da pesquisa (quando houver). Por exemplo: Agradecimentos à Fundação de Amparo à Pesquisa do Estado do Piauí (FAPEPI) pelo apoio financeiro referente ao Edital nº 001/2024. Processo nº 00110.000156/2024-69.

Todos os Autores devem atender a dois critérios mínimos:
Participar ativamente na construção e discussão dos resultados.
Revisar e aprovar a versão final do artigo.

A declaração de contribuições dos autores deverá constar em seção ao final do artigo para trabalhos escritos por 2 (dois) ou mais autores. A declaração deve incluir a contribuição individual de cada autor.

A partir do volume 10, ano de 2024, a revista adotará a Taxonomia de Papéis de Contribuidor (CRediT), o que significa que toda declaração de contribuição deve ser descrita usando as 14 (quatorze) funções de contribuidor do CRediT como segue: Conceitualização, Curadoria de dados, Análise Formal, Aquisição de financiamento, Investigação, Metodologia, Administração do projeto, Recursos, Software, Supervisão, Validação, Visualização, Redação - rascunho original e Redação - revisão e edição. As 14 (quatorze) funções estão disponíveis em \url{https://credit.niso.org/}.

Os autores deverão declarar se há algum conflito de interesses no momento da submissão do manuscrito. Caso não haja conflito de interesses, os autores deverão selecionar no formulário de submissão a opção: ``Os autores declaram não haver conflito de interesses''; caso contrário, deverá ser especificado em ``Comentários ao Editor''. Uma seção com a declaração de Conflito de Interesses deverá ser incluída após a seção de Declaração de autoria e contribuições dos autores.

Artigos envolvendo pesquisa com seres humanos é obrigatório a aprovação do Comitê de Ética em Pesquisa com Seres Humanos – CEP e apresentação do número do processo. Artigos que não contenham a documentação de aprovação do CEP serão automaticamente rejeitados.

As informações contidas nos tópicos acima devem ser devidamente preenchidas no como segue o quadro abaixo:

\quadroobrigatorio{Esta pesquisa teve o apoio financeiro da Fundação de Amparo à Pesquisa do Estado do Piauí (FAPEPI) pelo apoio financeiro referente ao Edital nº 001/2024. Processo nº 00110.000156/2024-69.}{Autor A: Conceituação; Autor B: Metodologia, Software, Curadoria de dados, Análise formal, Redação – rascunho original, Visualização, Investigação; Autor C: Redação – revisão e edição, Aquisição de financiamento, Supervisão. Todos os autores participaram ativamente da discussão dos resultados, revisaram e aprovaram a versão final do artigo.}{Os autores declaram que não têm interesse comercial ou associativo que represente um conflito de interesses em relação ao manuscrito.}{Esta pesquisa foi avaliada e aprovada pelo Comitê de Ética em Pesquisa com Seres Humanos – CEP do Instituto Federal do Piauí (Processo nº 081910).}

Ao final do texto deverão aparecer as \textbf{REFERÊNCIAS}, utilizando fonte \textbf{Arial Narrow}, tamanho da fonte 12, espaçamento simples entre linhas; com alinhamento à esquerda, espaçamento entre parágrafos em 0 pt antes e 0 pt depois. Uma referência deve ser separada da outra por apenas por espaço simples.

Todas as referências colocadas no artigo deverão seguir a NBR 6023:2018 da ABNT. Alguns exemplos seguem abaixo:

\textbf{Citação com citação} (Fonte Arial Narrow, 11, recuo 4 cm; obrigatório o número da página)

Freire (1989, p. 17):
\begin{citacao}
A relevância da biblioteca popular com relação aos programas de educação e de cultura popular em geral e não apenas de alfabetização de adultos, creio que é apreendida tanto por educadoras e educadores numa posição ingênua, ou astutamente ingênua, quanto por aquelas e aqueles que se inserem numa perspectiva crítica. O em que se distinguem, é na concepção - e na sua posta em prática - da biblioteca.
\end{citacao}

\textbf{Livros}:
SOBRENOME DO AUTOR, Prenomes abreviados. \textbf{Título do Livro}: subtítulo. Local de publicação: Editora, ano de publicação.

Exemplo:
FREIRE, P. \textbf{A importância do ato de ler}: em três artigos que se completam. São Paulo: Autores Associados: Cortez, 1989.

\textbf{Citação longa com mais de 3 autores} (Fonte Arial Narrow, 11, recuo 4 cm; obrigatório o número da página)

Os \textbf{artigos de revisão} são definidos por Noronha e Ferreira (2000, p. 191) como
\begin{citacao}
estudos que analisam a produção bibliográfica em determinada área temática, dentro de um recorte de tempo, fornecendo uma visão geral ou um relatório do estado-da-arte sobre um tópico específico, evidenciando novas ideias, métodos, subtemas que têm recebido maior ou menor ênfase na literatura selecionada (Noronha; Ferreira, 2000, p. 192).
\end{citacao}

\textbf{Capítulo de Livro}:
SOBRENOME DO AUTOR, Prenomes abreviados. Título do Capítulo: subtítulo. \textit{In}: SOBRENOME DO AUTOR, Prenomes abreviados (eds). \textbf{Título do Livro}. Local de publicação: Editora, ano de publicação. Páginas inicial e final (separados por hífen ``-'').

Exemplo:
NORONHA, D. P.; FERREIRA, S. M. S. P. Revisões de literatura. \textit{In}: CAMPELLO, B. S.; CONDÓN, B. V.; KREMER, J. M. (orgs.). \textbf{Fontes de informação para pesquisadores e profissionais}. Belo Horizonte: UFMG, 2000. p. 191-199.

ZANELLA, F. C. V.; MARTINS, C. F. Abelhas da Caatinga: biogeografia, ecologia e Conservação. \textit{In}: LEAL, I. R.; TABARELLI, M.; SILVA, J. M. C. (eds.). \textbf{Ecologia e Conservação da Caatinga}. Recife: Ed. Universitária da UFPE, 2003. p. 75-135.

\textbf{Citação de artigo} (Fonte Arial Narrow, 12; obrigatório o número da página)
Segundo Bocchiglieri, Mendonça e Henriques (2010, p. 175) ``a riqueza da mastofauna de médio e grande porte da Fazenda Jatobá é comparável à de algumas localidades mais estudadas e preservadas do Cerrado''.

\textbf{Artigo de Periódico}:
SOBRENOME DO AUTOR, Prenomes abreviados. Título do Artigo: subtítulo. \textbf{Título da Revista}, Número do volume, número do fascículo, páginas inicial e final do artigo (separados por hífen ``-'') e ano de publicação. DOI.

Exemplo:
BOCCHIGLIERI, A.; MENDONÇA, A. F.; HENRIQUES, R. P. B. Composição e diversidade de mamíferos de médio e grande porte no Cerrado do Brasil central. \textbf{Biota Neotropica}, v. 10, n. 3, p. 169-176, 2010. DOI: 10.1590/S1676-06032010000300019

\textbf{Citação de tese} (Fonte Arial Narrow, 12; obrigatório o número da página)
``O dia-a-dia da maioria dos professores está associado, por conta dos baixos salários, ao exercício de atividades em mais de uma escola e ao excesso de aulas'' (Aguiar, 2020, p. 198).

\textbf{Teses e Dissertações}:
SOBRENOME DO AUTOR, Prenomes com Abreviatura. \textbf{Título do trabalho}: subtítulo. Ano. Número de folhas. Dissertação ou Tese (Mestrado em ou Doutorado em) – Nome do Programa, Nome da Universidade, Local, Ano.

Exemplo:
AGUIAR, R. R. \textbf{A escola de tempo integral}: análise de uma proposta curricular inovadora para o ensino médio. 2020. 279 f. Tese (Doutorado em Ensino de Ciências) – Programa de Pós-Graduação Interunidades em Ensino de Ciências, Universidade de São Paulo-USP, São Paulo, 2020. Disponível em: https://www.teses.usp.br/teses/disponiveis/81/81131/tde-21122020-172638/en.php. Acesso em: 30 dez. 2023.

\textbf{Citação de página da internet} (Fonte Arial Narrow, 12)
``\textit{Trachycephalus} se enquadra filogeneticamente em Hylidae'' (Frost, 2023).

\textbf{Site/Homepage}:
SOBRENOME DO AUTOR, Prenomes abreviados. \textbf{Nome do site ou serviço}. Edição. Local: ano. Nº de pág. ou vol. (série) (se houver). Disponível em: endereço eletrônico. Acesso em: dia mês (abreviado), ano.

Exemplo:
FROST, D. R. 2023. \textbf{Amphibian Species of the World}: an Online Reference. Version 6.2 (Oct 2023). American Museum of Natural History, New York, USA, 2023. DOI 10.5531/db.vz.0001. Disponível em: https://amphibiansoftheworld.amnh.org/index.php.

- Solicitamos que as normas acima sejam cuidadosamente seguidas; caso contrário, os textos enviados \textbf{NÃO SERÃO RECEBIDOS} para avaliação. Em caso de dúvidas consulte a NBR 6023:2018 e NBR 10520:2023.

Atenção! Todas as citações deverão estar corretamente referenciadas.
